% OBJETIVOS

\section{Objetivos}

En este apartado se trata de definir los objetivos del trabajo(TFM). Hay que tener en cuenta que el documento a entregar como memoria de prácticas tiene que evidenciar la relación de competencias y contenidos específicos del título con el trabajo realizado. A continuación, se explican las directrices sobre cómo expresar los objetivos formalmente en la memoria del TFM.

Los objetivos en una memoria se definen como los logros que se quiere alcanzar durante el desarrollo del trabajo. La definición de los objetivos tiene que ser clara y concisa para no dar lugar a interpretaciones. Un ejemplo es comenzar con frases como: “Los objetivos establecidos para este proyecto…” “El objetivo principal es la implantación de un clasificador…”. Los objetivos tienen que cumplir que sean específicos, realistas y alcanzables.

Es posible que a lo largo del periodo de duración del trabajo los objetivos puedan sufrir alguna modificación. Ojo! Quizá también sea necesario cambiar el título del trabajo. Deberá consensuarlo con el tutor y comunicarlo al coordinador.

En la definición de objetivos, es conveniente separar entre generales y específicos. El/los objetivos generales se podrían considerar como la/las unidades más complejas para realizar el trabajo o proyecto. Estas unidades se deberían descomponer en otros objetivos más detallados o específicos, que son en metas más sencillas para poder alcanzar el objetivo general al que hace referencia.
